\begingroup
\section{Node TR: physical trace and relative Schur transfer}
\label{module:TR}
\noindent\textit{Audited source: \nolinkurl{TR_physical_trace.tex}.}
\subsection{The true straight-edge material trace}

Let $\alpha_i>0$, $\sum_i\alpha_i=2\pi$, and
$\theta_{i+1}-\theta_i=\alpha_i$.  Put
\begin{equation}\label{TR:eq:fan}
 a=\alpha_{i-1},\qquad b=\alpha_i,\qquad
 K_i=[\theta_i-a/2,\theta_i+b/2],\qquad x=\theta-\theta_i.
\end{equation}
The area-normalized tangential polygon has common support
\begin{equation}\label{TR:eq:h}
 h_\alpha=\left(\frac\pi{\sum_j\tan(\alpha_j/2)}\right)^{1/2}
\end{equation}
and radial boundary $h_\alpha\sec x$ on $K_i$.

For support increments $z=(z_i)$, the two endpoint displacement vectors
of side $i$ have components in the frame $(\nu_i,\tau_i)$
\begin{equation}\label{TR:eq:endpoint-components}
 (z_i,a_i^-),\qquad (z_i,a_i^+),
\end{equation}
where
\begin{equation}\label{TR:eq:apm}
 a_i^-=\frac{z_i-z_{i-1}}{\sin a}-z_i\tan\frac a2,
 \qquad
 a_i^+=\frac{z_{i+1}-z_i}{\sin b}+z_i\tan\frac b2.
\end{equation}
Indeed these are obtained by differentiating the two adjacent support-line
intersection equations.

The physical affine coordinate of the radial point
$h_\alpha\nu_i+h_\alpha\tan x\,\tau_i$ is
\begin{equation}\label{TR:eq:t}
 t_i(x)=\frac{\tan x+\tan(a/2)}
 {\tan(a/2)+\tan(b/2)}.
\end{equation}
Hence the genuine radial component of the straight-edge affine material
velocity is
\begin{equation}\label{TR:eq:Tphysical}
 (\widetilde{T}_\alpha z)_i(x)
 =z_i\cos x+
 \bigl((1-t_i(x))a_i^-+t_i(x)a_i^+\bigr)\sin x.
\end{equation}
For comparison, the stage-65 angular trace is
\begin{equation}\label{TR:eq:Tangular}
 (T_\alpha z)_i(x)
 =z_i\cos x+
 \bigl((1-s_i(x))a_i^-+s_i(x)a_i^+\bigr)\sin x,
 \qquad
 s_i(x)=\frac{x+a/2}{(a+b)/2}.
\end{equation}

We use the Fourier norm
\begin{equation}\label{TR:eq:Hhalf}
 \|f\|_{H^{1/2}}^2
 =\sum_{n\in\mathbb Z}(1+|n|)|\widehat f(n)|^2.
\end{equation}
All quotient estimates below are unchanged if the equivalent disk form
$q_{\mathbb D}$ is used.

\begin{lemma}[Endpoint-uniform interpolation jets]
\label{TR:lem:weight-jets}
Put $w=a+b$, $p=a/w$, $y=s_i(x)$, and assume $0<w\le w_0<\pi$.
Then, uniformly for $0\le p\le1$,
\begin{align}
 |t_i-s_i|&\le Cw^2s_i(1-s_i),
 \label{TR:eq:weight-value}\\
 w|t_i'-s_i'|&\le Cw^2.
 \label{TR:eq:weight-derivative}
\end{align}
In particular, no endpoint factor occurs in
\eqref{TR:eq:weight-derivative}.
\end{lemma}

\begin{proof}
In normalized variables,
\begin{equation}\label{TR:eq:normalized-t}
 \vartheta_{w,p}(y)=
 \frac{\tan(\frac w2(y-p))+\tan(\frac{wp}2)}
 {\tan(\frac{wp}2)+\tan(\frac{w(1-p)}2)},
 \qquad t_i(x)=\vartheta_{w,p}(y).
\end{equation}
The denominator is between $w/2$ and $C_{w_0}w$.  Direct
differentiation gives
\begin{equation}\label{TR:eq:normalized-second}
 \partial_y^2\vartheta_{w,p}(y)
 =\frac{w^2}{2}
 \frac{\sec^2(\frac w2(y-p))\tan(\frac w2(y-p))}
 {\tan(\frac{wp}2)+\tan(\frac{w(1-p)}2)},
\end{equation}
and hence $|\partial_y^2\vartheta_{w,p}|\le Cw^2$.
The function $f=\vartheta_{w,p}-y$ vanishes at $0$ and $1$.
The Green function of $-\partial_y^2$ on $[0,1]$ therefore gives
$|f(y)|\le Cw^2y(1-y)$.  Moreover
$\int_0^1f'(y)\,dy=0$, so
$|f'(y)|\le\|f''\|_\infty\le Cw^2$.  Since
$\partial_x=(2/w)\partial_y$, these are exactly
\eqref{TR:eq:weight-value}--\eqref{TR:eq:weight-derivative}.
\end{proof}

For the next lemma, write $w_i=\alpha_{i-1}+\alpha_i$ and
$d_i=a_i^+-a_i^-$.  The length of $K_i$ is $w_i/2$.

\begin{lemma}[Stable extraction of the internal quadratic coefficient]
\label{TR:lem:coefficient-stability}
For $S$ sufficiently small,
\begin{equation}\label{TR:eq:local-coefficient}
 w_i^2|d_i|^2
 \le C\left(
 [T_\alpha z]_{H^{1/2}(K_i)}^2+w_i^4|z_i|^2\right).
\end{equation}
On the exact area kernel, modulo exact support translations,
\begin{equation}\label{TR:eq:nodal-sampling}
 \sum_iw_i|z_i|^2
 \le C\|\PiNG T_\alpha z\|_{H^{1/2}(\T)}^2.
\end{equation}
Both constants are independent of adjacent ratios.
\end{lemma}

\begin{proof}
Pull $K_i$ to $0\le y\le1$ and put
\begin{equation*}
 Z=z_i,\qquad A=w_i a_i^-,\qquad B=w_i a_i^+,
 \qquad \xi=\frac{w_i}{2}(y-p).
\end{equation*}
The normalized trace is
\begin{equation}\label{TR:eq:normalized-trace}
 F_{w_i,p}(y)=Z\cos\xi+\{(1-y)A+yB\}\frac{\sin\xi}{w_i}.
\end{equation}
At $w_i=0$ its continuous limit is
\begin{equation}\label{TR:eq:homogeneous-normalized-trace}
 F_{0,p}(y)=Z+\frac12(y-p)\{(1-y)A+yB\}.
\end{equation}
Modulo constants, the last two normalized polynomials have coefficient
matrix
\begin{equation*}
 \begin{pmatrix}1+p&-p\\-1&1\end{pmatrix},
\end{equation*}
whose determinant is one and whose inverse is uniformly bounded for
$0\le p\le1$.  Finite-dimensional norm equivalence therefore gives
\begin{equation}\label{TR:eq:homogeneous-coercivity}
 |A|^2+|B|^2\le C[F_{0,p}]_{H^{1/2}(0,1)}^2.
\end{equation}
Taylor's formula in \eqref{TR:eq:normalized-trace}, in the same
finite-dimensional space, gives
\begin{equation}\label{TR:eq:normalized-perturbation}
 [F_{w_i,p}-F_{0,p}]_{H^{1/2}(0,1)}
 \le Cw_i^2(|Z|+|A|+|B|).
\end{equation}
After decreasing the mesh threshold, the $A,B$ terms are absorbed into
\eqref{TR:eq:homogeneous-coercivity}.  Since
$B-A=w_i d_i$ and the $H^{1/2}$ seminorm is invariant under affine
rescaling in one dimension, this proves \eqref{TR:eq:local-coefficient}.

We also need the constant and translation gauges.  The full normalized
$L^2$ three-jet matrix associated with \eqref{TR:eq:normalized-trace} is
uniformly invertible.  Hence, for every constant $c$,
\begin{equation}\label{TR:eq:local-sampling}
 \sum_iw_i|z_i-c|^2
 \le C\|T_\alpha z-c\|_{L^2(\T)}^2.
\end{equation}
Every first harmonic is exactly the angular trace of a support
translation $z_i^{\rm tr}=q\cdot\nu_i$, and $L_\alpha z^{\rm tr}=0$.
Subtract the unique translation producing the $n=\pm1$ part of
$T_\alpha z$.  For the resulting representative,
\begin{equation*}
 T_\alpha z=g+c,\qquad g=\PiNG T_\alpha z.
\end{equation*}
The exact area row is
\begin{equation}\label{TR:eq:exact-area-row}
 L_\alpha z=\sum_i\ell_i z_i=0,\qquad
 \ell_i=\tan\frac{\alpha_{i-1}}2+
        \tan\frac{\alpha_i}2\asymp w_i.
\end{equation}
By Cauchy--Schwarz and \eqref{TR:eq:local-sampling},
\begin{equation*}
 |c|L_\alpha\mathbf1=|L_\alpha(z-c\mathbf1)|
 \le C\Big(\sum_iw_i|z_i-c|^2\Big)^{1/2},
\end{equation*}
where $L_\alpha\mathbf1\asymp1$.  Thus
$|c|\le C\|g\|_{L^2}$.  Applying
\eqref{TR:eq:local-sampling} once more proves
\eqref{TR:eq:nodal-sampling}.  Exact translation subtraction does not
change $d_i$, because $a_i^-=a_i^+=q\cdot\tau_i$ for a translation.
\end{proof}

\begin{lemma}[Stable zero-wall bubble decomposition]
\label{TR:lem:zero-wall}
Let $E$ equal $E_i$ on $K_i$, where
\begin{equation}\label{TR:eq:trace-error-local}
 E_i=(t_i-s_i)d_i\sin x.
\end{equation}
Then
\begin{equation}\label{TR:eq:zero-wall-bound}
 \|E\|_{H^{1/2}(\T)}^2\le C\sum_iw_i^6|d_i|^2.
\end{equation}
\end{lemma}

\begin{proof}
Lemma~\ref{TR:lem:weight-jets} gives, in normalized coordinates,
\begin{equation}\label{TR:eq:E-normalized-jets}
 |E_i(y)|\le Cw_i^3|d_i|y(1-y),
 \qquad
 |\partial_yE_i(y)|\le Cw_i^3|d_i|.
\end{equation}
In particular, $E_i$ vanishes at both endpoints.  The same-cell
Gagliardo integral is at most $Cw_i^6|d_i|^2$.  For the cross-cell
integral use
\begin{equation*}
 |E_i(\theta)-E_j(\varphi)|^2
 \le2|E_i(\theta)|^2+2|E_j(\varphi)|^2.
\end{equation*}
For $\theta\in K_i$, integration of the circular kernel over
$\T\setminus K_i$ is bounded by
\begin{equation*}
 C\left(1+\operatorname {dist}(\theta,\partial K_i)^{-1}\right).
\end{equation*}
The endpoint factor in \eqref{TR:eq:E-normalized-jets} makes the resulting
Hardy integral at most $Cw_i^6|d_i|^2$.  Summing over $i$, and adding the
$L^2$ row, proves \eqref{TR:eq:zero-wall-bound}.  This argument sums the
whole complement of each cell at once, so it has no mesh-depth or
adjacent-ratio constant.
\end{proof}

\begin{lemma}[Physical versus angular interpolation]
\label{TR:lem:trace-close}
Let $S=\max_i\alpha_i+2\pi/N$.  On the exact area and translation
quotient,
\begin{equation}\label{TR:eq:trace-close}
 \|\PiNG(\widetilde T_\alpha-T_\alpha)z\|_{H^{1/2}(\T)}
 \le CS^2\|\PiNG T_\alpha z\|_{H^{1/2}(\T)}.
\end{equation}
The constant is independent of every adjacent-cell ratio.
\end{lemma}

\begin{proof}
Projection is contractive in \eqref{TR:eq:Hhalf}.  By
Lemmas~\ref{TR:lem:coefficient-stability} and \ref{TR:lem:zero-wall},
\begin{align*}
 \|\PiNG E\|_{H^{1/2}}^2
 &\le C\sum_iw_i^6|d_i|^2\\
 &\le CS^4\sum_i[T_\alpha z]_{H^{1/2}(K_i)}^2
   +C\sum_iw_i^8|z_i|^2\\
 &\le CS^4\|\PiNG T_\alpha z\|_{H^{1/2}}^2
   +CS^7\sum_iw_i|z_i|^2\\
 &\le CS^4\|\PiNG T_\alpha z\|_{H^{1/2}}^2.
\end{align*}
Here we use the representative constructed in the proof of
Lemma~\ref{TR:lem:coefficient-stability}, for which
$T_\alpha z=\PiNG T_\alpha z+c$; consequently its local seminorms are
those of $\PiNG T_\alpha z$.  Also $w_i\le2S$, and the same-cell
Gagliardo integrals are positive parts of the global seminorm.  Taking square roots proves
\eqref{TR:eq:trace-close}.
\end{proof}

\begin{remark}[The repaired terminology]
$T_\alpha z$ remains a useful trigonometric spline and the stage-65 disk
gap for it is valid.  Equation \eqref{TR:eq:Tphysical}, not
\eqref{TR:eq:Tangular}, is the trace of an actual affine map of a straight
edge.  The distinction is a small relative error, but it is not an exact
identity and may not be suppressed in a value comparison.
\end{remark}

\subsection{Transfer of D0 to the physical trace}

We first fix a normalization which was ambiguous in the historical
notation.  The reduced Fourier--Bessel form of F51 is denoted by
$\qDzero$; the physical half-Hessian form is
\begin{equation}\label{TR:eq:physical-disk-normalization}
 \boxed{\qphys(f,g):=2j_{0,1}^{\,2}\qDzero(f,g)
 =\frac12D^2\mathcal L(0)[fe_r,ge_r].}
\end{equation}
Every disk endpoint, metric, pairing, and distance below uses $\qphys$.
Thus the angular endpoint imported from F65 is first multiplied by
$2j_{0,1}^{\,2}$; this changes only its positive universal gap constant,
and changes neither its minimizer nor any relative trace comparison.

Let
\begin{equation}\label{TR:eq:source}
 b_\alpha=\PiNG\{h_\alpha\sec x-1\}
\end{equation}
on the cells $K_i$.  Define the angular and physical disk values
\begin{align}
 P_D(\alpha)&=\inf_{z\in Z_\alpha}
 \qphys(b_\alpha+\PiNG T_\alpha z),\label{TR:eq:PD}\\
 \widetilde P_D(\alpha)&=\inf_{z\in Z_\alpha}
 \qphys(b_\alpha+\PiNG\widetilde T_\alpha z).
 \label{TR:eq:PDphysical}
\end{align}
Here $Z_\alpha$ includes the exact area row and is quotiented by
translations.  Let
\begin{equation}\label{TR:eq:Gphysical}
 \widetilde\G_\alpha(z)
 =\qphys(\PiNG\widetilde T_\alpha z),
\end{equation}
and let $\widetilde z_\alpha^\sharp$ be the corresponding Schur corrector.

Write $H_{\rm ng}=\PiNG H^{1/2}(\T)$ and let
$\langle\cdot,\cdot\rangle_D$ be the bilinear form associated with
$\qphys$.  The Fourier--Bessel bounds from the disk Hessian give
\begin{equation}\label{TR:eq:Bessel-coercivity}
 c\|f\|_{H^{1/2}}^2\le \qphys(f)\le
 C\|f\|_{H^{1/2}}^2,\qquad f\in H_{\rm ng}.
\end{equation}
Set
\begin{equation}\label{TR:eq:trace-spaces}
 V_\alpha=\PiNG T_\alpha Z_\alpha,\qquad
 \widetilde V_\alpha=\PiNG\widetilde T_\alpha Z_\alpha.
\end{equation}
Thus $P_D(\alpha)$ and $\widetilde P_D(\alpha)$ are the squared
$\qphys$-distances from $-b_\alpha$ to $V_\alpha$ and
$\widetilde V_\alpha$, respectively.

\begin{lemma}[Differentiated Hilbert projection calculus]
\label{TR:lem:projection-calculus}
Let $q_t$ be uniformly coercive $C^2$ Hilbert forms on a fixed Hilbert
space, let $P_t,\widetilde P_t$ be the corresponding orthogonal
projections onto two $C^2$ families of closed subspaces, and put
$\Delta P_t=\widetilde P_t-P_t$.  Suppose, for $0\le r\le2$,
\begin{align}
 \|\partial_t^rq_t\|&\le C_rh^r,
 \label{TR:eq:q-material-bound}\\
 \|\partial_t^rP_t\|+\|\partial_t^r\widetilde P_t\|
 &\le C_rh^r,
 \qquad
 \|\partial_t^r\Delta P_t\|\le C_rh^r\varepsilon,
 \label{TR:eq:P-material-bound}\\
 \|\partial_t^rb_t\|&\le C_rh^rB.
 \label{TR:eq:b-material-bound}
\end{align}
Then
\begin{equation}\label{TR:eq:projection-calculus-output}
 \left|\partial_t^2\left\{
 \operatorname {dist}_{q_t}(-b_t,\widetilde V_t)^2-
 \operatorname {dist}_{q_t}(-b_t,V_t)^2\right\}\right|
 \le Ch^2\varepsilon B^2.
\end{equation}
\end{lemma}

\begin{proof}
Orthogonal completion of the square gives
\begin{equation}\label{TR:eq:projection-difference-identity}
 \operatorname {dist}_{q_t}(-b_t,\widetilde V_t)^2-
 \operatorname {dist}_{q_t}(-b_t,V_t)^2
 =-q_t(\Delta P_tb_t,(\widetilde P_t+P_t)b_t).
\end{equation}
Differentiate this identity twice.  Every resulting term contains one of
$\Delta P_t$, $\partial_t\Delta P_t$, or
$\partial_t^2\Delta P_t$; the remaining factors are bounded by
\eqref{TR:eq:q-material-bound}--\eqref{TR:eq:b-material-bound}.  Each term is
therefore at most $Ch^2\varepsilon B^2$.
\end{proof}

\begin{lemma}[Uniform material calculus for the paired trace spaces]
\label{TR:lem:material-calculus}
Let
\begin{equation}\label{TR:eq:quasiuniform-band}
 (1-\eta_0)a\le\gamma_i\le(1+\eta_0)a,\qquad
 \sum_i\gamma_i=2\pi,\qquad 0<\eta_0<\frac14,
\end{equation}
and let $\sum_i\delta_i=0$.  If
\begin{equation}\label{TR:eq:Rtr}
 R_{\rm tr}(\alpha)=\widetilde P_D(\alpha)-P_D(\alpha),
\end{equation}
then
\begin{equation}\label{TR:eq:transfer-hessian}
 |D^2R_{\rm tr}(\gamma)[\delta,\delta]|
 \le Ca\sum_i\gamma_i^2\delta_i^2.
\end{equation}
\end{lemma}

\begin{proof}
The assertion is trivial for $\delta=0$.  Otherwise put
\begin{equation}\label{TR:eq:material-h}
 h=\frac{\|\delta\|_{\ell^\infty}}a,
 \qquad \gamma(t)=\gamma+t\delta,
 \qquad |t|<c_0h^{-1}.
\end{equation}
Let $\Phi_t$ be the piecewise affine circle map carrying the normal nodes
of $\gamma$ to those of $\gamma(t)$, and pull all functions back by
$\Phi_t$.  For every oriented cyclic arc $I$ of length at most $\pi$,
\begin{equation}\label{TR:eq:arc-material}
 \left|t\sum_{i\in I}\delta_i\right|
 \le c_0a\# I\le Cc_0\sum_{i\in I}\gamma_i.
\end{equation}
Consequently $\Phi_t'$ and every chord ratio are a fixed small relative
perturbation of one.  Differentiating the pulled Gagliardo kernel
\begin{equation}\label{TR:eq:pulled-Gagliardo}
 \frac{\Phi_t'(\xi)\Phi_t'(\eta)}
 {\sin^2((\Phi_t(\xi)-\Phi_t(\eta))/2)}
\end{equation}
zero, one, or two times proves
\begin{equation}\label{TR:eq:form-material}
 \|\partial_t^rq_t\|\le C_rh^r,\qquad 0\le r\le2.
\end{equation}
Here the order-one part is exactly \eqref{TR:eq:pulled-Gagliardo}; the
$L^2$ row is elementary.  For the Bessel remainder use its established
order-minus-one coefficients
\begin{equation}\label{TR:eq:Bessel-differences}
 |\Delta^rk_n|\le C_r(1+|n|)^{-1-r},\qquad 0\le r\le3,
\end{equation}
and discrete summation by parts.  It satisfies the same estimate, with
one extra order of decay.

The normalized secant formula and two material differentiations give
\begin{equation}\label{TR:eq:source-material}
 \|\partial_t^rb_{\gamma(t)}\|_{H^{1/2}}
 \le C_rh^ra^{3/2},\qquad 0\le r\le2.
\end{equation}
Indeed a cell source has amplitude $O(a^2)$ and scale-invariant
$H^{1/2}$ charge $O(a^4)$; summing $O(a^{-1})$ cells gives $O(a^3)$.
Every normalized material derivative costs at most $h$.  The global
factor $h_{\gamma(t)}$ obeys the same bounds after differentiating its
explicit tangent sum.

We next fix the coefficient space.  Let
\[
 \tau_t^1=(\cos\theta_i(t))_i,\qquad
 \tau_t^2=(\sin\theta_i(t))_i
\]
be the two exact support-translation vectors.  They satisfy
$L_{\gamma(t)}\tau_t^\mu=0$.  With the area weights
$\ell_i(t)$ from \eqref{TR:eq:exact-area-row}, the $2\times2$ Gram matrix
\[
 M_t^{\mu\nu}=\sum_i\ell_i(t)\tau_{t,i}^\mu\tau_{t,i}^\nu
\]
is uniformly comparable with the identity in the quasiuniform band.
Define the explicit translation gauge
\begin{equation}\label{TR:eq:translation-gauge}
 \mathcal Q_tz=z-\sum_{\mu,\nu=1}^2
 (M_t^{-1})_{\mu\nu}
 \left(\sum_i\ell_i(t)z_i\tau_{t,i}^\nu\right)\tau_t^\mu
\end{equation}
and transport the area kernel by
\begin{equation}\label{TR:eq:area-transport}
 C_tz=\mathcal Q_t\left(
 z-\frac{L_{\gamma(t)}z}{L_{\gamma(t)}\mathbf1}\mathbf1\right).
\end{equation}
The coefficients of $L_{\gamma(t)}$ are comparable to $a$, and their
first two normalized derivatives are bounded by $C_rh^r$.  The same is
true of $\tau_t^\mu$, $M_t^{-1}$, $\mathcal Q_t$, and $C_t$.
The sampling estimate \eqref{TR:eq:nodal-sampling} makes the range of
$C_t$ uniformly transverse to the area and translation gauges.  Thus,
on the area/translation complement $X$ at $t=0$, the pulled trace injections
\begin{equation*}
 B_t=\PiNG(T_{\gamma(t)}C_t)\circ\Phi_t,
 \qquad
 \widetilde B_t=\PiNG(\widetilde T_{\gamma(t)}C_t)\circ\Phi_t
\end{equation*}
are uniformly bounded below in the pulled $q_t$ norm.

Apply the normalized calculation
\eqref{TR:eq:normalized-t}--\eqref{TR:eq:E-normalized-jets} on this
quasiuniform material band.  Its first two material derivatives give
\begin{align}
 \left|\partial_t^r\{(t_i-s_i)\sin x\}\right|
 &\le C_rh^ra^3y(1-y),\label{TR:eq:material-E-value}\\
 \left|\partial_y\partial_t^r\{(t_i-s_i)\sin x\}\right|
 &\le C_rh^ra^3,\qquad 0\le r\le2.
 \label{TR:eq:material-E-derivative}
\end{align}
The derivative row uses the bound without an endpoint factor from
\eqref{TR:eq:weight-derivative}.  The coefficient-extraction and zero-wall
proofs then apply after each material differentiation.  Together with
the derivative bounds for $C_t$, this gives
\begin{align}
 \|\partial_t^rB_t\|+\|\partial_t^r\widetilde B_t\|
 &\le C_rh^r,
 \label{TR:eq:B-material}\\
 \|\partial_t^r(\widetilde B_t-B_t)\|
 &\le C_rh^ra^2,\qquad 0\le r\le2.
 \label{TR:eq:DeltaB-material}
\end{align}
All operator norms use the fixed energy norm on $X$ and the pulled
$q_t$ norm on the target.  Notice that the endpoint factors were proved
before differentiating, so no inverse minimum cell is present.

For completeness, the $q_t$-orthogonal projection onto
$\operatorname {ran}B_t$ is
\begin{equation}\label{TR:eq:Gram-projection}
 P_t=B_tG_t^{-1}B_t^{*t},\qquad
 G_t=B_t^{*t}B_t,
\end{equation}
where $*t$ is the adjoint in $q_t$.  The corresponding formula with
$\widetilde B_t$ defines $\widetilde P_t$.  Uniform trace coercivity
bounds both Gram inverses.  Differentiating
\eqref{TR:eq:Gram-projection} and
\begin{equation*}
 \partial_tG_t^{-1}=-G_t^{-1}(\partial_tG_t)G_t^{-1}
\end{equation*}
twice, using \eqref{TR:eq:form-material}, \eqref{TR:eq:B-material}, and
\eqref{TR:eq:DeltaB-material}, yields
\begin{align}
 \|\partial_t^rP_t\|+\|\partial_t^r\widetilde P_t\|
 &\le C_rh^r,\notag\\
 \|\partial_t^r(\widetilde P_t-P_t)\|
 &\le C_rh^ra^2,\qquad 0\le r\le2.
 \label{TR:eq:projection-material}
\end{align}
Lemma~\ref{TR:lem:projection-calculus}, with
$\varepsilon=a^2$ and $B=a^{3/2}$, now gives
\begin{equation}\label{TR:eq:R-second-h}
 |D^2R_{\rm tr}(\gamma)[\delta,\delta]|
 \le Ch^2a^5=Ca^3\|\delta\|_{\ell^\infty}^2.
\end{equation}
Finally
$\sum_i\gamma_i^2\delta_i^2\ge ca^2
\|\delta\|_{\ell^\infty}^2$, which proves
\eqref{TR:eq:transfer-hessian}.
\end{proof}

\begin{lemma}[Relative physical-trace Schur transfer]
\label{TR:lem:physical-transfer}
There is a modulus $\omega_{\rm tr}(S)\to0$ such that
\begin{equation}\label{TR:eq:physical-transfer}
 |\{\widetilde P_D(\alpha)-P_D(\alpha)\}
  -\{\widetilde P_D(\alpha_*)-P_D(\alpha_*)\}|
 \le\omega_{\rm tr}(S)\Jfour(\alpha).
\end{equation}
Moreover $\widetilde\G_\alpha$ and $\G_\alpha$ are uniformly equivalent,
with relative error $O(S^2)$.
\end{lemma}

\begin{proof}
The equivalence of the support metrics follows from
Lemma~\ref{TR:lem:trace-close}, the Bessel coercivity
$\qphys\asymp H^{1/2}$, and polarization; more precisely,
\begin{equation}\label{TR:eq:metric-equivalence}
 |\widetilde\G_\alpha(z)-\G_\alpha(z)|
 \le CS^2\G_\alpha(z).
\end{equation}
Lemma~\ref{TR:lem:trace-close} also gives a graph bijection
$J_\alpha:V_\alpha\to\widetilde V_\alpha$ satisfying
$\|J_\alpha-I\|\le CS^2$ in the $\qphys$ norm.  The elementary Hilbert
distance comparison therefore gives
\begin{equation}\label{TR:eq:transfer-absolute}
 |R_{\rm tr}(\alpha)|
 \le CS^2\qphys(b_\alpha)\le CS^5.
\end{equation}
The final source estimate follows by normalized same/adjacent/separated
Gagliardo summation for the continuous secant source:
$\qphys(b_\alpha)\le C\sum_i\alpha_i^4\le CS^3$.

Put $a=2\pi/N$, $d_i=\alpha_i-a$, and recall the exact identity
\begin{equation}\label{TR:eq:Jfour-identity}
 \Jfour(\alpha)=
 \sum_i d_i^2(\alpha_i^2+2a\alpha_i+3a^2).
\end{equation}
If $\Jfour(\alpha)\ge S^{9/2}$, apply
\eqref{TR:eq:transfer-absolute} at $\alpha$ and at
$\alpha_*=a\mathbf1$ to obtain
\begin{equation}\label{TR:eq:far-transfer}
 |R_{\rm tr}(\alpha)-R_{\rm tr}(\alpha_*)|
 \le CS^{1/2}\Jfour(\alpha).
\end{equation}

Suppose instead that $\Jfour(\alpha)<S^{9/2}$.  Then
$\|d\|_{\ell^\infty}^4\le C\Jfour(\alpha)$, so
$\|d\|_{\ell^\infty}/S\le CS^{1/8}$.  Since
$S=a+\max_i\alpha_i$, the whole segment
$\gamma(t)=a\mathbf1+td$ lies in a fixed quasiuniform band and
$a\asymp S$.  Cyclic invariance makes
$DR_{\rm tr}(a\mathbf1)[d]=0$, because the gradient at the regular fan
is constant and $\sum_i d_i=0$.  Taylor's formula and
Lemma~\ref{TR:lem:material-calculus} give
\begin{align}
 |R_{\rm tr}(\alpha)-R_{\rm tr}(\alpha_*)|
 &\le Ca\int_0^1(1-t)
       \sum_i\gamma_i(t)^2d_i^2\,dt\notag\\
 &\le CS\Jfour(\alpha),
 \label{TR:eq:near-transfer}
\end{align}
where the last inequality follows from
\eqref{TR:eq:Jfour-identity} and quasiuniformity.  Combining
\eqref{TR:eq:far-transfer} and \eqref{TR:eq:near-transfer} proves
\eqref{TR:eq:physical-transfer} with
\begin{equation}\label{TR:eq:transfer-modulus}
 \omega_{\rm tr}(S)=C(S^{1/2}+S).
\end{equation}
\end{proof}

\begin{theorem}[Physical disk endpoint D0]
\label{TR:thm:physical-D0}
For all sufficiently fine positive fans,
\begin{equation}\label{TR:eq:physical-D0}
 \qphys(b_\alpha+\PiNG\widetilde T_\alpha z)
 -\widetilde P_D(\alpha_*)
 \ge c_f\Jfour(\alpha)
   +c_s\widetilde\G_\alpha(z-\widetilde z_\alpha^\sharp),
\end{equation}
where $c_f,c_s>0$ are independent of all minimum cells.
\end{theorem}

\begin{proof}
Complete the physical support square exactly:
\begin{equation*}
 \qphys(b_\alpha+\PiNG\widetilde T_\alpha z)
 =\widetilde P_D(\alpha)
  +\widetilde\G_\alpha(z-\widetilde z_\alpha^\sharp).
\end{equation*}
The angular endpoint gap of stages 65 and 129, together with
Lemma~\ref{TR:lem:physical-transfer}, gives
$\widetilde P_D(\alpha)-\widetilde P_D(\alpha_*)\ge c_f\Jfour$ after
decreasing the fine-mesh threshold.  Uniform coercivity gives the support
constant $c_s$.
\end{proof}

\subsection{An exact fixed-disk scalar for the polygon}

Before the harmless area dilation, let $p_i(h)$ be the intersection of
the support lines with normals $\nu_i,\nu_{i+1}$.  It depends linearly on
$(h_i,h_{i+1})$.  On side $i$ define
\begin{equation}\label{TR:eq:material-boundary}
 X_{\alpha,z}(\theta)
 =(1-t_i(x))p_{i-1}(h_\alpha\mathbf1+z)
  +t_i(x)p_i(h_\alpha\mathbf1+z).
\end{equation}
As $\theta$ traverses $K_i$, this traverses the actual straight side;
hence its image is exactly $\partial P(\alpha,h_\alpha\mathbf1+z)$.
Scale normalization of the eigenvalue makes the omitted final area
dilation irrelevant.

Write
\begin{equation}\label{TR:eq:W}
 X_{\alpha,z}(\theta)=e_r(\theta)+W_{\alpha,z}(\theta).
\end{equation}
At $z=0$, \eqref{TR:eq:material-boundary} is the radial parametrization of the
tangential polygon.  Equations \eqref{TR:eq:endpoint-components}--
\eqref{TR:eq:Tphysical} therefore give the exact identity
\begin{equation}\label{TR:eq:normal-W}
 \PiNG(W_{\alpha,z}\cdot e_r)
 =b_\alpha+\PiNG\widetilde T_\alpha z
 =:\widetilde\tau(\alpha,z).
\end{equation}

Choose a fixed annular extension of $W$ which equals $W$ at $r=1$ and
vanishes for $r\le1/2$.  For $\|W\|_{W^{1,\infty}}$ small, the map
$F_W={\rm id}+EW$ is bi-Lipschitz.  Pulling stiffness and mass back by
$F_W$ defines the scale-normalized first eigenvalue
\begin{equation}\label{TR:eq:LambdaW}
 \mathcal L(W)=\frac{|F_W(\mathbb D)|}{\pi}
 \lambda_1(F_W(\mathbb D)).
\end{equation}

\begin{lemma}[Parametric fixed-chart analyticity]
\label{TR:lem:parametric-analytic}
$\mathcal L$ is real analytic on a fixed small ball of
$W^{1,\infty}(\T;\mathbb R^2)$.  It depends only on the boundary image,
and
\begin{equation}\label{TR:eq:parametric-Hessian}
 D\mathcal L(0)=0,
 \qquad
 \frac12D^2\mathcal L(0)[W,W]
 =\qphys\bigl(\PiNG(W\cdot e_r)\bigr).
\end{equation}
\end{lemma}

\begin{proof}
The pullback coefficients are rational analytic functions of $DF_W$ and
$\det DF_W$ on a small $L^\infty$ ball.  Uniform ellipticity, simplicity
of the first eigenvalue, and the Riesz projection give operator-norm
analyticity exactly as in the radial fixed chart.  Two extensions with the
same boundary image are related by a bi-Lipschitz change of variables and
therefore give the same eigenvalue and area.  The disk is critical for the
scale-normalized functional.  Reparametrization invariance removes the
tangential component from its second differential; the normal component
is the Fourier--Bessel disk Hessian.  This proves
\eqref{TR:eq:parametric-Hessian}.
\end{proof}

Define the nonlinear spectral remainder
\begin{equation}\label{TR:eq:N}
 \mathcal N(W)=\mathcal L(W)-\mathcal L(0)
 -\qphys\bigl(\PiNG(W\cdot e_r)\bigr).
\end{equation}
Since \eqref{TR:eq:material-boundary} parametrizes the actual polygon,
\begin{equation}\label{TR:eq:exact-value-split}
 \boxed{
 \lambda_1(P(\alpha,H_\alpha(z)))
 =\mathcal L(0)+\qphys(\widetilde\tau(\alpha,z))
  +\mathcal N(W_{\alpha,z}).}
\end{equation}
Here the left side is area normalized; equivalently one may use the
unnormalized support polygon because \eqref{TR:eq:LambdaW} already performs
the scale normalization.
\endgroup
